\section{The exact larger-field profiles for an inducing $E_8$ block}
\label{sec:e8-inducing-fields}

\begin{theorem}[HD-E8IF1]
Let a nonsplit Weil sixfold source of exact imaginary quadratic
character field $K$ have period rank nine. Suppose its coefficient
representation is induced from an irreducible eight-dimensional
representation $\tau$ whose reflection normal closure is $W(E_8)$.
Assume that $K$ is one of the eight non-Eisenstein class-number-one
fields and that the inducing character field
$L=\Q(\chi_\tau)$ is strictly larger than $K$.
Then the complete list of necessary field and branch profiles is
\[
\begin{array}{c|c|c|c}
 K&L&\text{faithful inducing image }J&
     \text{real class of either outside factor}\\\hline
 \Q(i)\text{ or }\Q(\sqrt{-2})&\Q(\zeta_8)&
       \mu_8W(E_8)&D_8(a_3),\quad\Phi_8^2\\
 \Q(i)&\Q(\zeta_{12})&\mu_{12}W(E_8)&
       E_8(a_3),\quad\Phi_{12}^2.
\end{array}                                                    \tag{E8IF.1}
\]
In both rows the intermediate base is rational and has exactly
twelve branch points: ten real reflections and two outside inertias
of moving rank six. Write the outside elements as
$g_1=\lambda w_1$, $g_2=\varepsilon\lambda^{-1}w_2$,
where $w_i\in W(E_8)$ and $\varepsilon\in\{1,-1\}$.
In the first row $\lambda$ is primitive of order eight and each
outside spectrum is
\[
                  1^2,i^2,(-1)^2,(-i)^2.                    \tag{E8IF.2}
\]
In the second row $\lambda$ is primitive of order twelve and,
up to exchanging the two entries, their spectra are
\begin{align*}
 &(1,1,\zeta_6,\zeta_6,-1,-1,\zeta_3^2,\zeta_3^2),\\
 &(1,1,\zeta_6^{-1},\zeta_6^{-1},-1,-1,\zeta_3,\zeta_3).
\end{align*}
Equivalently the distinct exponent
fractions are $\{0,1/6,1/2,2/3\}$ and
$\{0,1/3,1/2,5/6\}$, each repeated twice.
These necessary profiles have Chevalley--Weil signature $(3,3)$.
They do not assert the existence of a generating tuple, a dominant
component, or a nonsplit actual polarization.
\end{theorem}

\begin{proof}
The normalizer of the real $E_8$ reflection group is scalar times
$W(E_8)$: reduce to a real orthogonal normalizer by Cohen's
Lemma~2.12, then use the chamber argument and the absence of
a diagram automorphism, as in D-BI-13. The group already contains
$-I$ and is defined over $\Q$. Consequently
$J\subset\C^*W(E_8)$ and $W(E_8)\subset J$.
Every finite-order element of $J$ is $\lambda w$ with $w\in W(E_8)$
and $\lambda$ a root of unity. The scalar is determined modulo
$\{1,-1\}$, and this class defines a character of $J$.
Indeed the ambiguity is exactly the scalar center of $W(E_8)$.

Frobenius reciprocity identifies the original complex multiplicity
variation with the actual intermediate-cover variation for $\tau$.
Its Euler equation is $6=16(h-1)+D$. Genus at least two is impossible;
at genus one the marked-base dimension is at most six. Thus rank nine
forces genus zero and a genuine product-one generating inertia tuple
with
\[
             r\ge12,\qquad\sum_{j=1}^r\operatorname{rank}(1-g_j)=22.
                                                               \tag{E8IF.4}
\]
Let $t$ count the inertias outside $W(E_8)$. Such an inertia has a
nonreal scalar phase. Since $w$ is real orthogonal, a nonreal
eigenvalue and its conjugate have equal multiplicities, at most four.
Each outside inertia therefore has moving rank at least four.
Every other genuine inertia has moving rank at least one, so
$22\ge r+3t$, and hence $t\le3$.
The scalar product-one relation excludes $t=1$; $t=0$ would give
$J=W(E_8)$ and a real character field, contrary to $K\subset L$.

Suppose $t=3$. The three outside moving ranks sum to at most thirteen;
each is at least four, so each is at most five. Each corresponding
$w_j$ has $\lambda_j^{-1}$ as an eigenvalue with multiplicity at
least three. Since $w_j$ is rational, the minimal polynomial of
$\lambda_j$ occurs at least three times in its degree-eight
characteristic polynomial. Thus $[\Q(\lambda_j):\Q]\le2$.
The nonreal phases have order three, four or six, and their classes
modulo $\{1,-1\}$ have order three or two. In the product of three
nonidentity such classes, the number of order-two entries must be even.
Two order-two entries cancel but leave a nonidentity order-three entry;
zero order-two entries is the only possibility. All three phases
therefore have order three or six. The group and all character values
are defined over $\Q(\zeta_3)$, contradicting the non-Eisenstein
quadratic subfield $K$. This excludes $t=3$.

We have $t=2$. The two scalar phases multiply to $\varepsilon=\pm1$,
and their moving ranks have sum at most twelve. At least one fixes a
space of dimension at least two. Its rational real factor shows that
$[\Q(\lambda):\Q]\le4$. Because the tuple generates $J$ and
$W(E_8)\subset J$, its scalar matrices generate the same cyclotomic
field as $\lambda$. Their traces are $8\lambda^a$, so
\[
                         L=\Q(\lambda).                    \tag{E8IF.5}
\]
The orders of a nonreal root of unity of degree at most four are
$3,4,5,6,8,10,12$. Orders three, four and six give a quadratic
field and cannot satisfy $K\subsetneq L$. Orders five and ten
give $\Q(\zeta_5)$, whose only quadratic subfield is real.
The three quadratic subfields of $\Q(\zeta_8)$ have radicands
$2,-1,-2$, and those of $\Q(\zeta_{12})$ have radicands
$3,-1,-3$. This gives exactly the two field rows in (E8IF.1).
These elementary subfield calculations are also certified in D-BI-11.

In either remaining row the phase has degree four. On either real
factor its multiplicity as an eigenvalue is at most two, so both
outside moving ranks are at least six. Equation (E8IF.4) now forces
$r=12$, both ranks equal to six, and all ten other entries to be
reflections. Each real factor has characteristic polynomial
$\Phi_d^2$, where $d=8$ or $12$; all eight dimensions are exhausted
by the four Galois-conjugate eigenvalues, each with multiplicity two.
Since $J$ contains $W(E_8)$ and the primitive scalar $\lambda I$,
it is the asserted full scalar extension.

For $d=8$, adding a primitive eighth-root exponent to
$\{1,3,5,7\}/8$ yields $\{0,1,2,3\}/4$; each occurs twice.
The moving rank is six and the age is three. For $d=12$, the
primitive exponents are $\{1,5,7,11\}/12$. Adding the scalar
exponent gives one of the two displayed exponent sets; the reciprocal phase,
even after adding $1/2$ for $\varepsilon=-1$, gives the other set.
The ages are $8/3$ and $10/3$. Thus in either row the two outside
ages sum to six, while ten reflections contribute five. Subtracting
the coefficient degree eight in Chevalley--Weil gives three, and
the conjugate calculation gives three as well.

For the class names in (E8IF.1), the complete $E_8$ class list is
the published Carter Table~11, pp.~52--54, with the documented
CHEVIE class representatives already checked in D-GAUSS
\cite{WG:Carter,WG:CHEVIE}. The present exact certificate reconstructs
the two word matrices from the retained rational Cartan matrix,
checks their orders and characteristic polynomials, and verifies that
they are the only classes in that complete list with $\Phi_8^2$
or $\Phi_{12}^2$. The class sizes are respectively $3\,628\,800$
and $2\,419\,200$. No product-one existence assertion is inferred
from this local class calculation.
\end{proof}
